% Standalone article. Compile twice with pdflatex.
\documentclass[11pt,a4paper]{article}
\usepackage[T1]{fontenc}
\usepackage{lmodern,amsmath,amssymb,booktabs,microtype}
\usepackage[a4paper,top=22mm,bottom=22mm,left=25mm,right=25mm]{geometry}
\usepackage[hidelinks]{hyperref}
\hypersetup{pdftitle={DimTypes and Boost.Units for orbital mechanics}}
\setlength{\parindent}{0pt}
\setlength{\parskip}{5pt plus 1pt minus 1pt}
\setlength{\emergencystretch}{1.3em}
\newcommand{\code}[1]{\texttt{#1}}
\newcommand{\Fp}{\mathbb{F}_p}
\newcommand{\heading}[1]{\par\medskip{\large\bfseries #1}\par\nobreak}
\begin{document}

{\LARGE\bfseries DimTypes and Boost.Units for orbital mechanics\par}
\smallskip
{\large Finite-field exponents, unit conversion, and numerical functions\par}
\medskip

\textbf{Abstract.}
An ephemeris may supply a gravitational parameter in
$\mathrm{AU}^3/\mathrm{day}^2$ while an orbit integrator works in kilometres
and seconds. Both DimTypes and Boost.Units can check the dimensions of the
calculation. They differ in which unit expressions they admit and how they
represent exponents. DimTypes fits SpaceBallistics because one declaration
defines the project's unit vocabulary, conversions are visible in the
equations, and the same quantity types reach into the numerical code.
Boost.Units has the stronger model when independently defined systems,
mixed-unit products, or exact distinctions among many rational powers matter.

\heading{1\quad Two algebras for the same physical equation}
For an elliptic orbit with semimajor axis $a$ and gravitational parameter
$\mu$, the characteristic time is $\tau=\sqrt{a^3/\mu}$. Since
$[a]=L$ and $[\mu]=L^3T^{-2}$, the type of $\tau$ should be a time.
Writing $\sqrt{a^2/\mu}$ leaves an unwanted $L^{-1/2}$.
Either library can expose this error when the result is assigned to a time.
Omitting the factor $2\pi$ in the period $P=2\pi\tau$ is a physics error
that passes either dimensional check.

Boost.Units represents a dimension as a canonical list of base-dimension
types and reduced \code{static\_rational} exponents. A unit pairs a dimension
with a system, and a \code{quantity<Unit,Y>} carries the value of type
\code{Y}. Algebraically equivalent expressions reduce to the same dimension.
Distinct rational powers stay distinct within the range of the library's
integer arithmetic~\cite{boostdim,boostunits}.

DimTypes represents a quantity as \code{DimQ<E,U,RepT,MaxDims>}. The two
64-bit template arguments pack the exponent vector $E$ and the unit identifiers
$U$ into fixed-width fields. Only the magnitude of type \code{RepT} is stored
in the object. Unit identifiers are ordinary indices. Each \emph{exponent} is
a residue in the prime field $\Fp=\mathbb{Z}/p\mathbb{Z}$. For $p\nmid b$,
the encoding is
\[
  \rho(a/b)=(a\bmod p)\,(b\bmod p)^{-1}\pmod p .
\]
Multiplication adds exponent fields, division subtracts them, and taking a
square root multiplies each field by $2^{-1}$. With the default $p=251$,
$2^{-1}=126$. Squaring a square root restores the original exponent because
$2\cdot126\equiv1\pmod{251}$. The magnitude still uses ordinary numerical
arithmetic~\cite{dimtypes}.

\heading{2\quad What fits in one exponent field?}
Two reduced fractions $a/b$ and $c/d$ have the same residue precisely when
$p$ divides $ad-bc$, assuming both denominators are invertible. Thus
$L^{126}$ and $L^{1/2}$ have the same encoded dimension at $p=251$.
The implementation computes the largest $H$ for which all reduced fractions
of height $|a|+b\leq H$ ($b>0$) are distinguishable. The table gives the
computed \code{MaxHeight} values. Unit conversion and formatting recover a
fraction by searching this range. Arithmetic itself works on residues, so an
out-of-range exponent can silently acquire another exponent's type or fail
when decoding is required~\cite{dimtypes}.

\begin{center}\small
\begin{tabular}{@{}lrrr@{}}\toprule
Maximum base dimensions & 7 & 8 (default) & 9\\
Bits per exponent & 9 & 8 & 7\\
Prime modulus $p$ & 509 & 251 & 127\\
Collision-free height $H$ & 24 & 18 & 13\\
\bottomrule\end{tabular}
\end{center}

\newpage
\heading{3\quad Declaring the units and using them}
The following DimTypes vocabulary suffices for a two-body calculation.
Each row names a dimension, its fundamental unit, and additional units with
their scales in that fundamental unit. The declaration and calculation below
belong in the same application namespace~\cite{dimtypes}.

\begin{verbatim}
DECLARE_DIMS(double, 8,
  (Len,  m,   (km, 1000.0), (AU, 149597870700.0)),
  (Time, sec, (day, 86400.0)))

auto mu_eph = 2.959122e-4 * Cube(1.0_AU) / Sqr(1.0_day);
auto mu     = To_Time_sec(To_Len_km(mu_eph));
auto a      = 149597870.7_km;
auto tau    = SqRt(Cube(a) / mu);
static_assert(std::is_same_v<decltype(tau), Time_sec>);
\end{verbatim}

The solar parameter is rounded for this example. The macro generates the
aliases (including \code{Len\_km} and the representation-polymorphic
\code{Len\_km\_T<R>}), literal suffixes, unit names and scales, conversion
functions, and output support. The two conversions of $\mu$ operate on a
\emph{derived} quantity. The length conversion cubes the AU-to-km scale,
and the time conversion divides by $86400^2$. The result is about
$1.32712\times10^{11}\,\mathrm{km}^3/\mathrm{s}^2$, giving
$\tau\approx5.02264\times10^6\,\mathrm{s}$. A maintainer can read the
unit policy at the point where ephemeris data enter the calculation.

The equivalent form $1/\tau=\sqrt{\mu}/a^{3/2}$ carries a
length exponent $3/2$ of height five until it cancels. Ordinary \code{constexpr} integer functions
compute the result types. For symbolic manipulations that can generate large
intermediate powers even when the answer is simple, Boost's rational
representation gives a stronger guarantee. Neither design needs to store the
exponent vector in each numerical value. The choice alone establishes no
runtime or compilation-speed advantage.

For a calculation already written in a supplied system, Boost can require
less setup. The equivalent of the example's square root is
\code{root<2>(pow<3>(a) / mu)} for compatible Boost quantities.
In SpaceBallistics the same small vocabulary appears in declarations,
constants, diagnostics, and conversion sites. Generating those from one list
saves the work of maintaining corresponding application aliases and helpers.

There is a cost to the declaration's brevity. The macro assigns dimension
and unit slots by declaration order. The application namespace owns the
generated names, but its identity is absent from \code{DimQ}'s four template
arguments. Independent declarations that assign different meanings to the
same slots can therefore produce identical C++ types. SpaceBallistics benefits
from sharing one declaration across components. Boost's named base-dimension
and base-unit types are better suited to separately authored vocabularies.

\newpage
\heading{4\quad Mixed units, numerical functions, and application fit}
In DimTypes, \code{1.0\_m + 1.0\_km} requires an explicit conversion.
So does \code{1.0\_m * 1.0\_km}. Both operands occupy the length slot,
which can hold just one unit. The compound $\mathrm{km}/\mathrm{s}$ uses
two different slots and is directly representable. Boost supports
\emph{heterogeneous} products. Its documentation constructs an area from
metres times centimetres and converts the result to square metres. It also
supplies SI and CGS systems and distinguishes absolute temperatures from
temperature differences, including offset conversions. A DimTypes
\code{DECLARE\_DIMS} scale is multiplicative and cannot by itself express a
Celsius origin~\cite{boostunits,boostexamples,dimtypes}.

\textbf{Where CEMaths matters.}
The type check on $\sqrt{a^3/\mu}$ determines a result dimension.
Computing its magnitude is a separate numerical problem. The distinction
matters elsewhere in an orbit solver too. The quantity overload for
\code{ATan2(y,x)} accepts two real coordinates with matching dimensions and
compatible units, then calls the scalar two-argument function. The quantity
overload for \code{Exp} admits a dimensionless argument. DimTypes couples
the quantity overloads to CEMaths and allows registered scalar
representations, including an application-supplied automatic-differentiation
type with the requisite operations. Boost also supports custom value types
and quantity-aware mathematical functions. DimTypes gives the project control
over the scalar algorithms and their exceptional cases. For example, its
custom \code{Exp} computes the reduced argument using a split $\ln 2$.
For double-precision $x=650$, that change reduced the measured relative error
from about $204$ to $0.67$ machine epsilons. Such control entails testing
accuracy, special values, and toolchain-dependent constant evaluation. On
the tested Clang~23 build, \code{std::round} and \code{std::ldexp} prevent
constant evaluation of the general \code{Exp} path~\cite{dimtypes,boostexamples}.

\textbf{Choice for SpaceBallistics.}
DimTypes fits the existing orbit code because its exponents are small, its
conversions identify where unit conventions change, and its scalar functions
are already shared with the numerical code. Coordinate frame, body identity,
and time scale are additional distinctions supplied by SpaceBallistics above
the scalar layer. A length in an Earth-fixed frame has the same scalar
dimension as a length in an inertial frame~\cite{merkin}. Boost.Units is the
more natural choice when interfaces must preserve heterogeneous units or
combine independent unit systems, or when a finite-field collision would
invalidate the intended dimensional guarantee.

\begingroup\small
\begin{thebibliography}{5}\setlength{\itemsep}{1pt}\setlength{\parskip}{0pt}
\bibitem{dimtypes} L.~Merkin. \href{https://github.com/LMerkin/DimTypes/tree/233c44e95650cc575bd5f9fe99b2256281aa3965}{\emph{DimTypes}, revision \texttt{233c44e}}.
See \code{Encodings.hpp}, \code{DeclareDims.h}, \code{DimTypes.hpp},
\code{CEMaths.hpp}, and \code{Tests/DimTest.cpp}.
\bibitem{boostdim} M.~C.~Schabel and S.~Watanabe. \emph{Boost.Units},
\href{https://www.boost.org/doc/libs/latest/doc/html/boost_units/Dimensional_Analysis.html}{Dimensional Analysis}.
\bibitem{boostunits} M.~C.~Schabel and S.~Watanabe. \emph{Boost.Units},
\href{https://www.boost.org/doc/libs/latest/doc/html/boost_units/Units.html}{Units} and
\href{https://www.boost.org/doc/libs/latest/doc/html/boost_units/Quantities.html}{Quantities}.
\bibitem{boostexamples} M.~C.~Schabel and S.~Watanabe. \emph{Boost.Units},
\href{https://www.boost.org/doc/libs/latest/doc/html/boost_units/Examples.html}{Examples} and
\href{https://www.boost.org/doc/libs/latest/doc/html/boost_units/Reference.html}{Library Reference}.
\bibitem{merkin} L.~Merkin.
\href{https://meetups.yadro.com/static/documents/template_metaprogramming_in_modern_cpp_for_space_ballistics_problems_merkin.pdf}{\emph{Template metaprogramming in modern C++ for space ballistics}}.
Presentation slides (in Russian).
\end{thebibliography}
\endgroup
\end{document}
